\chapter{The Stability of Superposed Fluids: The Kelvin--Helmholtz Instability}
\label{ch:11}

%% SKELETON -- sections to be filled incrementally

\section{The perturbation equations}\label{sec:100}

As we have already remarked in \S\,90, the Kelvin--Helmholtz instability arises when we consider the character of the equilibrium of a stratified heterogeneous fluid when the different layers are in relative motion. The initial stationary state whose stability we wish to examine, then, is that of an incompressible, inviscid fluid in which there is horizontal streaming. To be specific, we shall suppose that the streaming takes place in the $x$-direction with a velocity $U$. The assumption that the fluid is inviscid allows us to consider $U$ as an arbitrary function of the height $z$.

Let the actual density at any point $(x,y,z)$, as the result of a disturbance, be $\rho+\delta\rho$. Let the corresponding change in the pressure be $\delta p$, and finally, let the components of the velocity in the perturbed state be $U+u$, $v$, and $w$. The equations governing the perturbation are (cf.\ equations X~(8), (9), and (28)):
\begin{equation}
\label{eq:11-1}
\rho\frac{\partial u}{\partial t}+\rho U\frac{\partial u}{\partial x}+\rho w\frac{dU}{dz} = -\frac{\partial}{\partial x}\,\delta p,
\end{equation}
\begin{equation}
\label{eq:11-2}
\rho\frac{\partial v}{\partial t}+\rho U\frac{\partial v}{\partial x} = -\frac{\partial}{\partial y}\,\delta p,
\end{equation}
\begin{equation}
\label{eq:11-3}
\rho\frac{\partial w}{\partial t}+\rho U\frac{\partial w}{\partial x} = -\frac{\partial}{\partial z}\,\delta p - g\,\delta\rho + \sum_s T_s\!\left(\frac{\partial^2}{\partial x^2}+\frac{\partial^2}{\partial y^2}\right)\!\delta z_s\,\delta(z-z_s),
\end{equation}
\begin{equation}
\label{eq:11-4}
\frac{\partial}{\partial t}\,\delta\rho + U\frac{\partial}{\partial x}\,\delta\rho = -w\frac{d\rho}{dz},
\end{equation}
\begin{equation}
\label{eq:11-5}
\frac{\partial}{\partial z}\,\delta z_s + U_s\frac{\partial}{\partial x}\,\delta z_s = w(z_s),
\end{equation}
and
\begin{equation}
\label{eq:11-6}
\frac{\partial u}{\partial x}+\frac{\partial v}{\partial y}+\frac{\partial w}{\partial z} = 0.
\end{equation}

In writing equation~\eqref{eq:11-3} we have allowed for the possibility that at some prescribed levels $z_s$ the density may change discontinuously and bring into play effects due to surface tension. Further, in equation~\eqref{eq:11-5} a subscript $s$ distinguishes the value of the quantity at $z = z_s$.

Analysing the disturbance into normal modes, we seek solutions whose dependence on $x$, $y$, and $t$ is given by\footnote{Notice that we are now assuming a time dependence $e^{nt}$, in contrast to Chapter~X where it was $e^{nt}$.}
\begin{equation}
\label{eq:11-7}
\exp\,\mathrm{i}(k_x\,x+k_y\,y+nt).
\end{equation}
For solutions having this dependence on $x$, $y$, and $t$, equations~\eqref{eq:11-1}--\eqref{eq:11-6} become:
\begin{equation}
\label{eq:11-8}
\mathrm{i}\rho(n+k_x U)u + \rho(DU)w = -\mathrm{i}k_x\,\delta p,
\end{equation}
\begin{equation}
\label{eq:11-9}
\mathrm{i}\rho(n+k_x U)v = -\mathrm{i}k_y\,\delta p,
\end{equation}
\begin{equation}
\label{eq:11-10}
\mathrm{i}\rho(n+k_x U)w = -D\delta p - g\,\delta\rho - k^2\sum_s T_s\,\delta z_s\,\delta(z-z_s),
\end{equation}
\begin{equation}
\label{eq:11-11}
\mathrm{i}(n+k_x U)\delta\rho = -wD\rho,
\end{equation}
\begin{equation}
\label{eq:11-12}
\mathrm{i}(n+k_x U_s)\delta z_s = w_s,
\end{equation}
and
\begin{equation}
\label{eq:11-13}
\mathrm{i}(k_x u+k_y v) = -Dw,
\end{equation}
where $D$ denotes $d/dz$.

Multiplying equations~\eqref{eq:11-8} and~\eqref{eq:11-9} by $-\mathrm{i}k_x$ and $-\mathrm{i}k_y$, respectively, adding, and making use of equation~\eqref{eq:11-13}, we obtain
\begin{equation}
\label{eq:11-14}
\mathrm{i}\rho(n+k_x U)Dw - \mathrm{i}\rho k_x(DU)w = -k^2\delta p;
\end{equation}
whereas, by combining equations~\eqref{eq:11-10}, \eqref{eq:11-11}, and~\eqref{eq:11-12}, we have
\begin{equation}
\label{eq:11-15}
\mathrm{i}\rho(n+k_x U)w = -D\delta p - \mathrm{i}\rho(Dp)\frac{w}{n+k_x U} + \mathrm{i}k^2\sum_s T_s\!\left(\frac{w}{n+k_x U_s}\right)_{\!s}\!\delta(z-z_s).
\end{equation}

Now eliminating $\delta p$ between equations~\eqref{eq:11-14} and~\eqref{eq:11-15}, we obtain
\begin{multline}
\label{eq:11-16}
D\!\left[\rho(n+k_x U)Dw - \rho k_x(DU)w\right] - k^2\rho(n+k_x U)w \\
= gk^2\bigg[D\rho - \frac{k^2}{g}\sum_s T_s\,\delta(z-z_s)\bigg]\frac{w}{n+k_x U}.
\end{multline}

If we suppose that the fluid is confined between two rigid planes at $z=0$ and $z=d$ (say), then we must require that the solutions of equation~\eqref{eq:11-16} must satisfy the boundary conditions
\begin{equation}
\label{eq:11-17}
w = 0 \quad\text{at } z = 0 \quad\text{and}\quad z = d.
\end{equation}

The fact that equation~\eqref{eq:11-16} involves a $\delta$-function means that to interpret the equation correctly at the interface $z = z_s$ (for example) we must integrate the equation over an infinitesimal element $(z_s-\epsilon, z_s+\epsilon)$ and pass to the limit $\epsilon = 0$. Before we effect such an integration and carry out the stated limiting process, it is necessary to observe that
\begin{equation}
\label{eq:11-18}
w/(n+k_x U) \quad\text{must be continuous at an interface.}
\end{equation}
If $U$ is continuous at an interface, \eqref{eq:11-18} simply reduces to a condition requiring the continuity of $w$; but if $U$ should be discontinuous at $z = z_s$, then we must require, instead, the uniqueness of the normal displacement of any point on the interface; and according to equation~\eqref{eq:11-12} this latter condition is equivalent to~\eqref{eq:11-18}.

Now integrating equation~\eqref{eq:11-16} between $z_s-\epsilon$ and $z_s+\epsilon$ and passing to the limit $\epsilon = 0$, we obtain, in view of~\eqref{eq:11-18},
\begin{equation}
\label{eq:11-19}
\Delta_s\!\left[\rho(n+k_x U)Dw - \rho k_x(DU)w\right] = gk^2\!\left[\Delta_s(\rho)-\frac{k^2T_s}{g}\right]\!\left(\frac{w}{n+k_x U}\right)_{\!s}\!,
\end{equation}
where
\begin{equation}
\label{eq:11-20}
\Delta_s(f) = f_{s+\epsilon,0} - f_{s-\epsilon,0}
\end{equation}
is the jump a quantity $f$ experiences at $z = z_s$, and the subscript $a$ denotes the value which a quantity, known to be continuous at an interface, takes at the interface $z = z_s$. We may consider equation~\eqref{eq:11-19} as a boundary condition which must be satisfied at a surface of discontinuity, while the equation
\begin{equation}
\label{eq:11-21}
D\!\left[\rho(n+k_x U)Dw - \rho k_x(DU)w\right] - k^2\rho(n+k_x U)w = gk^2(D\rho)\frac{w}{n+k_x U}
\end{equation}
is valid everywhere else.

On expanding equation~\eqref{eq:11-21} and rearranging, we have
\begin{multline}
\label{eq:11-22}
(n+k_x U)(D^2-k^2)w - k_x(D^2U)w - g\,k^2\frac{D\rho}{\rho}\frac{w}{n+k_x U} \\
+ \frac{D\rho}{\rho}\!\left[(n+k_x U)Dw - k_x(DU)w\right] = 0.
\end{multline}
It may be noted here that the last term in equation~\eqref{eq:11-22}, appearing with the factor $D\rho/\rho$, represents the effect of the heterogeneity of the fluid on the inertia; in most cases of interest, the neglect of this effect, in comparison with the effect on the potential energy, is justified.

\section{The case of two uniform fluids in relative horizontal motion separated by a horizontal boundary}\label{sec:101}

Let two uniform fluids of densities $\rho_1$ and $\rho_2$ be separated by a horizontal boundary at $z = 0$. Let the density $\rho_2$ of the upper fluid be less than the density $\rho_1$ of the lower fluid so that, in the absence of streaming, the arrangement is a stable one. We shall further suppose that the two fluids are streaming with the constant velocities $U_1$ and $U_2$.

In each of the two regions of constant $\rho$ and $U$, equation~\eqref{eq:11-21} reduces to
\begin{equation}
\label{eq:11-23}
(D^2-k^2)w = 0.
\end{equation}
(A factor $n+k_x U$ has been discarded; this is permissible under the circumstances.) The general solution of equation~\eqref{eq:11-23} is a linear combination of the integrals $e^{+kz}$ and $e^{-kz}$. Since $w$ cannot increase exponentially on either side of the interface at $z = 0$ and since $w/(n+k_x U)$ must be continuous on this surface, we can write
\begin{equation}
\label{eq:11-24}
w_1 = A(n+k_x U_1)e^{+kz} \quad (z < 0)
\end{equation}
and
\begin{equation}
\label{eq:11-25}
w_2 = A(n+k_x U_2)e^{-kz} \quad (z > 0)
\end{equation}
as the solutions appropriate for the two regions. Applying the condition~\eqref{eq:11-19} to the solution given by these equations, we obtain the characteristic equation
\begin{equation}
\label{eq:11-26}
\rho_1(n+k_x U_1)^2 + \rho_2(n+k_x U_2)^2 = gk\!\left(\rho_1-\rho_2\right) + \frac{k^2 T}{g}\;.
\end{equation}

Letting
\begin{equation}
\label{eq:11-27}
\alpha_1 = \frac{\rho_1}{\rho_1+\rho_2} \quad\text{and}\quad \alpha_2 = \frac{\rho_2}{\rho_1+\rho_2},
\end{equation}
we can rewrite equation~\eqref{eq:11-26} in the form
\begin{equation}
\label{eq:11-28}
\alpha_1(n+k_x U_1)^2+\alpha_2(n+k_x U_2)^2 = gk\!\left(\alpha_1-\alpha_2\right)+\frac{k^2 T}{g(\rho_1+\rho_2)}.
\end{equation}

Expanding equation~\eqref{eq:11-28}, we have
\begin{equation}
\label{eq:11-29}
n^2+2k_x(\alpha_1 U_1+\alpha_2 U_2)n+k_x^2(\alpha_1 U_1^2+\alpha_2 U_2^2) - gk\!\left(\alpha_1-\alpha_2\right)+\frac{k^2T}{g(\rho_1+\rho_2)} = 0.
\end{equation}

The roots of this equation are given by
\begin{multline}
\label{eq:11-30}
n = -k_x(\alpha_1 U_1+\alpha_2 U_2) \pm \bigg[gk\!\left(\alpha_1-\alpha_2\right)+\frac{k^2 T}{g(\rho_1+\rho_2)} \\
{}-k_x^2\alpha_1\alpha_2(U_1-U_2)^2\bigg]^{1/2}\!.
\end{multline}

\subsection*{(a) The case when surface tension is absent}

In the absence of surface tension, equation~\eqref{eq:11-30} gives
\begin{equation}
\label{eq:11-31}
n = -k_x(\alpha_1 U_1+\alpha_2 U_2)\pm\!\left[gk(\alpha_1-\alpha_2)-k_x^2\alpha_1\alpha_2(U_1-U_2)^2\right]^{1/2}\!.
\end{equation}

Several conclusions of interest can be drawn from this equation.

(1) When $k_x = 0$,
\begin{equation}
\label{eq:11-32}
n = \pm\left[gk(\alpha_1-\alpha_2)\right]^{1/2},
\end{equation}
i.e.\ perturbations transverse to the direction of streaming are unaffected by its presence.

(2) In every other direction, instability occurs when
\begin{equation}
\label{eq:11-33}
k_x^2\alpha_1\alpha_2(U_1-U_2)^2 > gk(\alpha_1-\alpha_2);
\end{equation}
i.e.\ for a given difference in velocity $U_1-U_2$ and for a given direction of the wave vector $k$, instability occurs for all wave numbers
\begin{equation}
\label{eq:11-34}
k > \frac{g(\alpha_1-\alpha_2)}{\alpha_1\alpha_2(U_1-U_2)^2\cos^2\theta},
\end{equation}
where $\theta$ is the angle between the directions of $\mathbf{k}$ and $\mathbf{U}$. According to~\eqref{eq:11-34}, for a given $(U_1-U_2)$, instability occurs for the least wave number when $\mathbf{k}$ is in the direction of $\mathbf{U}$; and this minimum wave number, $k_{\min}$, is given by
\begin{equation}
\label{eq:11-35}
k_{\min} = \frac{g(\alpha_1-\alpha_2)}{\alpha_1\alpha_2(U_1-U_2)^2}.
\end{equation}
We have instability for $k > k_{\min}$.

The striking aspect of the instability predicted by~\eqref{eq:11-34} and~\eqref{eq:11-35} is that it occurs no matter how small $U_1-U_2$ may be. The stability of this static arrangement, in the absence of streaming, is unable to inhibit the instability, in the presence of streaming, for disturbances of sufficiently small wavelengths. This is the Kelvin--Helmholtz instability. Helmholtz stated this result in the following form: `\textit{Every perfect geometrically sharp edge by which a fluid flows must tear it asunder and establish a surface of separation, however slowly the rest of the fluid may move.}'

One may say that the Kelvin--Helmholtz instability arises by the crinkling of the interface by the shear which is present; and this crinkling occurs even for the smallest differences in the velocities of the two fluids. The question occurs whether this latter circumstance may not, in part, be due to the discontinuous distributions of $\rho$ and $U$ which have been considered; and may not be present if the distributions of $\rho$ and $U$ are continuous and the gradient of $U$ is sufficiently small compared to the gradient of $\rho$. We return to this question in \S\S\,102--4.

\subsection*{(b) The stabilizing effect of surface tension}

When surface tension is present, equation~\eqref{eq:11-30} will predict stability if
\begin{equation}
\label{eq:11-36}
k_x^2\alpha_1\alpha_2(U_1-U_2)^2 < gk\!\left(\alpha_1-\alpha_2\right)+\frac{k^2T}{g(\rho_1+\rho_2)}.
\end{equation}
(We have set $k_x = k$ in equation~\eqref{eq:11-30}, since these are the disturbances which are most sensitive to the Kelvin--Helmholtz instability.) An alternative form of~\eqref{eq:11-36} is
\begin{equation}
\label{eq:11-37}
\alpha_1\alpha_2(U_1-U_2)^2 < g\!\left(\frac{\alpha_1-\alpha_2}{k}+\frac{kT}{g(\rho_1+\rho_2)}\right)\!.
\end{equation}
The right-hand side of this inequality has a minimum when
\begin{equation}
\label{eq:11-38}
\frac{\alpha_1-\alpha_2}{k^2} = \frac{T}{g(\rho_1+\rho_2)}.
\end{equation}
If $k_*$ denotes the value of $k$ given by~\eqref{eq:11-38}, we shall have stability if
\begin{equation}
\label{eq:11-39}
(U_1-U_2)^2 < \frac{2g}{k_*}\frac{\alpha_1-\alpha_2}{\alpha_1\alpha_2}.
\end{equation}

Inserting for $k_*$ in accordance with~\eqref{eq:11-38}, we conclude that \textit{surface tension will suppress the Kelvin--Helmholtz instability if}
\begin{equation}
\label{eq:11-40}
(U_1-U_2)^2 < \frac{2}{\alpha_1\alpha_2}\sqrt{\frac{Tg(\alpha_1-\alpha_2)}{\rho_1+\rho_2}}.
\end{equation}
This result is due to Kelvin.

For air over sea-water ($\alpha_2 = 0.00126$; $\rho_1 = 1.02$ g/cm$^3$; $T = 74$ dyne/cm; and $g = 981$ cm/sec$^2$)~\eqref{eq:11-40} predicts that we will have stability for
\begin{equation}
\label{eq:11-41}
|U_1-U_2| < 650\text{ cm/sec};
\end{equation}
and when $|U_1-U_2|$ has this maximum value compatible with stability, we find
\begin{equation}
\label{eq:11-42}
k_* = 3{\cdot}68\text{ cm}^{-1}, \quad \lambda_* = 2\pi/k_* = 1{\cdot}71\text{ cm},
\end{equation}
\[
n_* = \alpha_2 k_*|U_1-U_2| = 3{\cdot}02\text{ sec}^{-1},
\]
and
\[
n/k_* = 0{\cdot}82\text{ cm/sec.}
\]

Therefore, when $|U_1-U_2|$ exceeds 650 cm/sec (12.5 nautical miles per hour) instability will manifest itself as surface waves with wavelength 1.71 cm and wave velocity 0.82 cm/sec. Commenting on this result, Kelvin stated in his original paper: `Observation shows the sea to be ruffled by wind of much smaller velocity than this. Such ruffling is, therefore, due to the viscosity of the air.' The first part of this statement of Kelvin has generally been repeated though other causes have been advanced for the `discrepancy'. However, Munk has recently pointed out that a wind velocity of 650 cm/sec is associated with several phenomena observed on the surface of the seas. The number of breaking waves (`white caps') suddenly increases; the convection processes, as indicated by the soaring of sea gulls, seem to change; and the rate of evaporation also increases suddenly. Munk ascribes all these to a sudden change in the pattern of air-flow over the waves consequent to the onset of the Kelvin--Helmholtz instability at the predicted wind velocity of 650 cm/sec.

In this connexion we may also refer to some experiments of Francis in which air was blown over lubricating oil in a wind-tunnel. Francis observed that at a sharply defined critical wind-speed small ripples appeared which were unstable and grew rapidly (see Fig.\ 117). The critical wind-speeds observed in this manner are well in accord with those predicted by~\eqref{eq:11-40}. It would thus appear that these experiments demonstrate, at any rate, the occurrence of Kelvin--Helmholtz instability under controlled laboratory conditions.

\figurePlaceholder{117}{Photograph by J.\ R.\ D.\ Francis of an unstable wave produced by air blown over a viscous oil in a wind tunnel; in the illustration, the wave has just grown from small ripples. Wind from left to right; camera looking slightly upwind and downward on to oil surface. Ripples show as transverse bars of light.}

\section{The effect of a continuous variation of $U$ on the development of the Kelvin--Helmholtz instability}\label{sec:102}

As we have remarked, the striking aspect of the occurrence of the Kelvin--Helmholtz instability at the interface between two fluids, in the absence of surface tension, is its independence on the magnitude of $|U_1-U_2|$. It is a matter of some interest to determine whether this aspect of the Kelvin--Helmholtz instability is due to the sharp discontinuities in $\rho$ and $U$ which has been assumed in its derivation. Before we consider this matter generally in \S\S\,103 and 104, we shall examine the case of two superposed fluids of different densities, separated by a transition layer of intermediate density in which the velocity of streaming varies continuously from that of the lower fluid to that of the upper fluid. Specifically, the distribution of $\rho$ and $U$ which we shall investigate is the following (see Fig.\ 118):
\begin{equation}
\label{eq:11-43}
\begin{aligned}
z > +d, &\quad \rho = \rho_0(1-\epsilon), &\quad U &= U_0, \\
+d > z > -d, &\quad \rho = \rho_0, &\quad U &= U_0 z/d, \\
z < -d, &\quad \rho = \rho_0(1+\epsilon), &\quad U &= -U_0.
\end{aligned}
\end{equation}

\figurePlaceholder{118}{The distributions of $\rho$ and $U$ for which the Kelvin--Helmholtz instability is investigated.}

Since $U$ is continuous at the interfaces at $z = +d$ and $z = -d$, the boundary conditions which must be satisfied at these surfaces are (cf.\ equations~\eqref{eq:11-18} and~\eqref{eq:11-19})
\begin{equation}
\label{eq:11-44}
w \text{ is continuous at } z = +d \text{ and } -d
\end{equation}
and
\begin{equation}
\label{eq:11-45}
\Delta_{\pm d}\!\left[\rho(n+k_x U)(DU)w\right] - gk^2\Delta_{\pm d}(\rho)w_{\pm d} = 0.
\end{equation}

The equation satisfied by $w$ in each of the regions of constant $\rho$ is the same equation~\eqref{eq:11-23}. We may accordingly write for the solutions in the three regions:
\begin{equation}
\label{eq:11-46}
\begin{aligned}
w_+ &= A_+\,e^{-kz} && (z > d), \\
w_0 &= A_0\,e^{-kz}+B_0\,e^{+kz} && (d > z > -d), \\
w_- &= B_-\,e^{+kz} && (z < -d).
\end{aligned}
\end{equation}
The solutions in the regions $z > d$ and $z < d$ have been so chosen that they vanish at $+\infty$ and $-\infty$, respectively. The continuity of $w$ at $z = +d$ and $-d$ leads to the conditions
\begin{equation}
\label{eq:11-47}
\begin{aligned}
A_+\,e^{-kd} &= A_0\,e^{-kd}+B_0\,e^{+kd} = w_{+d}, \\
B_-\,e^{-kd} &= A_0\,e^{+kd}+B_0\,e^{-kd} = w_{-d}.
\end{aligned}
\end{equation}

We shall restrict our further discussion to perturbations for which $k_y = 0$ and $k_x = k$ since instability will set in the $x$-direction for the least wave number. By applying the condition~\eqref{eq:11-45} at $z = \pm d$ to the solutions given in~\eqref{eq:11-46}, we obtain (in the case $k_y = 0$):
\begin{equation}
\label{eq:11-48}
(1-\epsilon)\!\left[(n+kU_0)+g\epsilon k(A_0\,e^{-kd}+B_0\,e^{+kd})\right] + \left[(n+kU_0)U_0/d+gk\right]\!(A_0\,e^{-kd}+B_0\,e^{+kd}) = 0
\end{equation}
and
\begin{equation}
\label{eq:11-49}
(1+\epsilon)\!\left[-(n-kU_0)-B_0\,e^{-kd}\right] - \left[(n-kU_0)U_0/d + gk\right](A_0\,e^{+kd}+B_0\,e^{-kd}) = 0.
\end{equation}

These equations can be reduced to give
\begin{equation}
\label{eq:11-50}
-\frac{A_0}{B_0} = \frac{d\!\left[(n+kU_0)^2+g\epsilon k\right]+(n+kU_0)U_0/d - 2(n+kU_0)^2/d}{d\!\left[(n+kU_0)^2+g\epsilon k\right](A_0/B_0)},
\end{equation}
and
\begin{equation}
\label{eq:11-51}
-\frac{B_0}{A_0} = \frac{d\!\left[(n-kU_0)^2-g\epsilon k\right]+(n-kU_0)U_0/d + 2(n-kU_0)^2/d}{d\!\left[(n-kU_0)^2-g\epsilon k\right](B_0/A_0)}.
\end{equation}

Let
\begin{equation}
\label{eq:11-52}
\nu = n/kU_0 \quad\text{and}\quad \kappa = 2kd,
\end{equation}
so that $\nu$ and $\kappa$ measure $n$ and $k$ in the units $U_0 k$ and $1/(2d)$, respectively. Further, let
\begin{equation}
\label{eq:11-53}
J = \frac{\epsilon g d}{U_0^2} = \frac{1}{2}\frac{g\epsilon}{\rho_0 (U/2d)^2};
\end{equation}
this non-dimensional `number' has the meaning
\begin{equation}
\label{eq:11-54}
J = -\frac{g}{2}\frac{\Delta\rho/2d}{\rho_0(U/2d)^2},
\end{equation}
where $\Delta\rho$ ($= -2\epsilon\rho_0$) is the difference in the density of the two fluids between which, in the transition layer, the velocity changes continuously from $-U$ (in the lower fluid) to $+U$ (upper fluid); it is the Richardson number for the problem on hand. (For the general definition of the Richardson number, its meaning, and its significance, see \S\S\,103 and 104 below.)

In terms of $\nu$, $\kappa$, and $J$, equations~\eqref{eq:11-50} and~\eqref{eq:11-51} take the more convenient forms
\begin{equation}
\label{eq:11-55}
\frac{A_0}{B_0}\,e^{-2\kappa} = 1 - \frac{\kappa(\nu+1)^2}{J+(\nu+1)+\frac{1}{2}\epsilon\kappa(\nu+1)^2}
\end{equation}
and
\begin{equation}
\label{eq:11-56}
\frac{B_0}{A_0}\,e^{-2\kappa} = 1 - \frac{\kappa(\nu-1)^2}{J-(\nu-1)+\frac{1}{2}\epsilon\kappa(\nu-1)^2}.
\end{equation}

Eliminating $A_0/B_0$, we obtain the characteristic equation
\begin{equation}
\label{eq:11-57}
e^{-2\kappa} = \left[1-\frac{\kappa(\nu+1)^2}{J+(\nu+1)+\frac{1}{2}\epsilon\kappa(\nu+1)^2}\right]\!\left[1-\frac{\kappa(\nu-1)^2}{J-(\nu-1)+\frac{1}{2}\epsilon\kappa(\nu-1)^2}\right]\!.
\end{equation}

While there is no formal difficulty in discussing equation~\eqref{eq:11-57} quite generally, the main interest in the present problem lies in cases in which the change in density is so small that we may neglect it in all terms except in the combination $J$. This neglect is equivalent to omitting the last term in equation~\eqref{eq:11-22} which occurs with the factor $D\rho/\rho$; as we have already remarked, its omission is equivalent to ignoring the effect of the heterogeneity of the fluid on the inertia while retaining its effect on the potential energy. On this approximation, equation~\eqref{eq:11-57} gives
\begin{equation}
\label{eq:11-58}
\left[(J+1)^2-\nu^2\right]e^{-2\kappa} = \left[(J+\nu+1)-\kappa(\nu+1)^2\right]\!\left[(J-\nu+1)-\kappa(\nu-1)^2\right]\!.
\end{equation}

On expansion and simplification, equation~\eqref{eq:11-58} becomes
\begin{equation}
\label{eq:11-59}
\kappa^2\nu^4-(2J\kappa-2\kappa^2-2\kappa+1-e^{-2\kappa})\nu^2+\left[(J+1)(1-e^{-2\kappa})-\kappa\right]\!\left[(J+1)(1+e^{-2\kappa})-\kappa\right] = 0.
\end{equation}

For stability, $\nu^2$ should clearly be real and positive. First, we may verify that the roots $\nu^2$ of equation~\eqref{eq:11-59} are always real. By evaluating the discriminant $\Delta$ of the equation, we find
\begin{equation}
\label{eq:11-60}
\Delta = 4\kappa^2(J+1)^2 e^{-2\kappa}+f(4\kappa J+1-e^{-2\kappa}),
\end{equation}
where
\begin{equation}
\label{eq:11-61}
f = (2\kappa-1)^2-e^{-2\kappa} = (2\kappa-1-e^{-\kappa})(2\kappa-1+e^{-\kappa}).
\end{equation}
The expression for $\Delta$ can be rewritten in the form
\begin{equation}
\label{eq:11-62}
\Delta = 4\kappa^2(J+1)^2 e^{-2\kappa}+4\kappa(J+1)f+f(-4\kappa^2 f)
\end{equation}
or
\begin{equation}
\label{eq:11-63}
\Delta = \left[2\kappa(J+1)e^{-\kappa}+e^{\kappa} f\right]^2-e^{2\kappa} f^2+f(-4\kappa^2 f).
\end{equation}
The last two terms in equation~\eqref{eq:11-63} can be combined to give
\begin{equation}
\label{eq:11-64}
\Delta = \left[2\kappa(J+1)e^{-\kappa}+e^{\kappa} f\right]^2-f\!\left[(2\kappa-1)e^{+\kappa}+e^{-\kappa}\right]^2.
\end{equation}

Now let $\kappa_0$ be such that
\begin{equation}
\label{eq:11-65}
2\kappa_0-1 = e^{-\kappa_0} \quad (\kappa_0 = 0{\cdot}7388).
\end{equation}
Then (cf.\ equation~\eqref{eq:11-61})
\begin{equation}
\label{eq:11-66}
f > 0 \quad\text{for } \kappa > \kappa_0 \quad\text{and}\quad f < 0 \quad\text{for } \kappa < \kappa_0.
\end{equation}
since $2\kappa > (1-e^{-\kappa})$ for all $\kappa > 0$. Hence, one of the two roots must be negative if
\begin{equation}
\label{eq:11-68}
\frac{\kappa}{1+e^{-\kappa}} < J+1 < \frac{\kappa}{1-e^{-\kappa}}.
\end{equation}

\figurePlaceholder{119}{The region of instability in the $(J,\kappa)$-plane for the distributions shown in Fig.\ 118; $\kappa$ is the wave number in the unit $1/(2d)$ and $J$ ($=\epsilon g d/U_0^2$) is the Richardson number.}

For values of the Richardson number in the range specified in~\eqref{eq:11-68}, the motion is unstable; and outside this range, it is stable.\footnote{Notice that when $J = 0$ the motion is stable for $\kappa > 1{\cdot}2785$ ($1{\cdot}2785$ is the root of the equation $\kappa = 1+e^{-\kappa}$); this result is due to Rayleigh.} The region of instability in the $(J,\kappa)$-plane is shown in Fig.\ 119. It is seen that for any assigned value of the Richardson number, there is a band of wavelengths for which the motion is unstable. For large $J$ the critical wave number, in the immediate vicinity of which there is instability, approaches $\kappa = J+1$; the corresponding wave velocity approaches the shear velocity ($= 2U_0$) at the transition layer.

Goldstein has extended these considerations to the case when the transition layer is further subdivided so that the density changes from $\rho_0+\epsilon$ to $\rho_0-\epsilon$ in five steps
\[
(\rho_0+\epsilon \to \rho_0+\tfrac{1}{2}\epsilon \to \rho_0+\rho_0-\tfrac{1}{2}\epsilon \to \rho_0-\epsilon)
\]
instead of three, while the velocity increases linearly through the three intermediate layers. Goldstein found that in this case, also, there are bands of wavelengths for which the motion is unstable for any assigned value of the Richardson number. Moreover, it appears that by further subdivision of the transition layer we do not gain anything: there are always bands of wavelengths for which the motion is unstable and we cannot avoid the Kelvin--Helmholtz instability for any value of the Richardson number.

\section{The Kelvin--Helmholtz instability in a fluid in which both $\rho$ and $U$ are continuously variable}\label{sec:103}

We shall preface the consideration of the mathematical aspects of the problem by seeking the origin of the Kelvin--Helmholtz instability in the manner we sought the origin of Rayleigh's criterion for the stability of revolving fluids in \S\,66.

The source of the Kelvin--Helmholtz instability clearly lies in the kinetic energy stored in the relative motion of the different layers: the tendency towards mixing and instability will be greater, the larger the prevailing shear as measured by $dU/dz$. The only counteracting force damping this tendency is derived from inertia; and so long as inertia can maintain a sufficient pressure gradient and prevent the mixing, instability will not occur. To obtain a quantitative expression of this condition, suppose that two equal neighbouring volumes, at heights $z$ and $z+\delta z$, are interchanged. The work that must be done to effect this interchange against the acceleration of gravity, per unit volume, is given by
\begin{equation}
\label{eq:11-69}
\delta W = -g\,\delta\rho\,\delta z,
\end{equation}
where $\delta\rho$ is the difference in the density at the two heights. The kinetic energy which is available to do this work (again, per unit volume) is given by
\begin{equation}
\label{eq:11-70}
\tfrac{1}{2}\rho\!\left[U^2+(U+\delta U)^2-\tfrac{1}{2}(U+U+\delta U)^2\right] = \tfrac{1}{4}\rho(\delta U)^2.
\end{equation}
A sufficient condition for stability is, therefore,
\begin{equation}
\label{eq:11-71}
\tfrac{1}{4}\rho(\delta U)^2 < -g\,\delta\rho\,\delta z.
\end{equation}
An equivalent form of this condition is
\begin{equation}
\label{eq:11-72}
\left(\frac{dU}{dz}\right)^{\!2} < -\frac{4}{g}\frac{d\rho}{\rho\;dz}.
\end{equation}
The number (cf.\ equation~\eqref{eq:11-54})
\begin{equation}
\label{eq:11-73}
J = -\frac{g}{\rho}\frac{d\rho/dz}{(dU/dz)^2}
\end{equation}
which measures the ratio of the buoyancy force to the inertia is the general definition of the Richardson number. The inequality~\eqref{eq:11-72}, therefore, suggests that a \textit{necessary condition for stability is that the Richardson number be everywhere greater than $\tfrac{1}{4}$}:
\begin{equation}
\label{eq:11-74}
J \geqslant \tfrac{1}{4} \quad\text{for stability.}
\end{equation}

The arguments leading to~\eqref{eq:11-74} do not, of course, justify one's inferring that its violation must lead to instability: for, it can very well happen that although stored energy is available for initiating instability, no mechanism exists for transforming the energy into possible hydrodynamic modes. The mathematical difficulty of establishing a general condition such as~\eqref{eq:11-74}, under suitable restrictions, appears to stem from this cause.

\subsection*{(a) Analytical results for the case $\rho = \rho_0 e^{-\beta z}$ and $U = U_0 z/d$}

Mathematically, the simplest distributions for which one may seek to establish a condition such as~\eqref{eq:11-74} are
\begin{equation}
\label{eq:11-75}
\rho = \rho_0\,e^{-\beta z} \quad\text{and}\quad U = U_0 z/d \quad (z > 0),
\end{equation}
where $\beta$ and $d$ are constants. The Richardson number is then a constant throughout the fluid and has the value
\begin{equation}
\label{eq:11-76}
J = g\beta d^2/U_0^2.
\end{equation}

For the distributions of $\rho$ and $U$ given in~\eqref{eq:11-75}, equation~\eqref{eq:11-22} becomes
\begin{equation}
\label{eq:11-77}
(n+kU_0 z/d)^2\!\left[D^2w - \beta Dw\right]+\left[-k^2+\frac{\beta g k^2}{(n+kU_0 z/d)^2}\right]w = 0.
\end{equation}

By measuring $z$ in the unit $d$ and letting
\begin{equation}
\label{eq:11-78}
n = U_0\nu/d, \quad k = \kappa/d, \quad\text{and}\quad \beta = \lambda/d,
\end{equation}
we can rewrite equation~\eqref{eq:11-77} in the form
\begin{equation}
\label{eq:11-79}
(\nu+\kappa z)^2\!\left[D^2w - \lambda Dw\right]+\left[-\kappa^2+\frac{\kappa^2 J}{(\nu+\kappa z)^2}\right]w = 0.
\end{equation}

Solutions of this equation must be sought which satisfy the boundary conditions
\begin{equation}
\label{eq:11-80}
w \to 0 \quad\text{for } z = 0 \quad\text{and}\quad z \to \infty.
\end{equation}

By the change of variable
\begin{equation}
\label{eq:11-81}
w = e^{\lambda z}W,
\end{equation}
equation~\eqref{eq:11-79} reduces to
\begin{equation}
\label{eq:11-82}
(\nu+\kappa z)^2\frac{d^2W}{dz^2}+\left[-(\kappa^2+\lambda^2)+\frac{\kappa^2 J}{(\nu+\kappa z)^2}\right]W = 0.
\end{equation}

Now let
\begin{equation}
\label{eq:11-83}
\zeta = (\kappa+\nu/\kappa)\sqrt{4\kappa^2+\lambda^2},
\end{equation}
and
\begin{equation}
\label{eq:11-84}
j = \sqrt{\tfrac{1}{4}(4\kappa^2+\lambda^2)} \quad\text{and}\quad m^2 = \tfrac{1}{4}-J.
\end{equation}

Equation~\eqref{eq:11-82} then simplifies to the form
\begin{equation}
\label{eq:11-85}
\frac{\zeta^2 d^2W}{d\zeta^2}+\left(-\tfrac{1}{4}+\frac{\tfrac{1}{4}-m^2}{\zeta^2}\right)W = 0.
\end{equation}

If we were to discard the factor $\zeta^2$ in equation~\eqref{eq:11-84}, then
\begin{equation}
\label{eq:11-86}
\mathscr{L}W = \frac{d^2W}{d\zeta^2}+\left(-\frac{1}{4}+\frac{\frac{1}{4}-m^2}{\zeta^2}\right)W = 0.
\end{equation}

We recognize in equation~\eqref{eq:11-86} Whittaker's standard form of the equation for the confluent hypergeometric function. Indeed, the condition at infinity requires that the solution of equation~\eqref{eq:11-86}, appropriate to the problem on hand, be, in fact, Whittaker's function,\footnote{For the definition and properties of Whittaker's function see E.\ T.\ Whittaker and G.\ N.\ Watson, \textit{Modern Analysis} (Cambridge: at the University Press, 1927), chapter XVI, particularly pp.\ 339--42.} $W_{j,m}(\zeta)$. Thus
\begin{equation}
\label{eq:11-87}
W = W_{j,m}(\zeta).
\end{equation}
This solution has the asymptotic behaviour
\begin{equation}
\label{eq:11-88}
W_{j,m}(\zeta) \sim \zeta^j e^{-\zeta/2} \quad (\zeta \to \infty).
\end{equation}

The corresponding form of the solution for $w$ is
\begin{equation}
\label{eq:11-89}
w(\zeta) = e^{\lambda z/2}W_{j,m}\!\left[(\kappa+\nu/\kappa)\sqrt{4\kappa^2+\lambda^2}\right]\!.
\end{equation}

The boundary condition at $z = 0$ now requires
\begin{equation}
\label{eq:11-90}
W_{j,m}\!\left[\frac{\nu}{\kappa}\sqrt{4\kappa^2+\lambda^2}\right] = 0.
\end{equation}
Hence,
\begin{equation}
\label{eq:11-91}
\frac{\nu}{\kappa}\sqrt{4\kappa^2+\lambda^2} \quad\text{must be a zero of } W_{j,m}(\zeta).
\end{equation}

Now Dyson has investigated the zeros of the Whittaker function $W_{j,m}(x)$ and has shown that it has no complex zeros; that it has an infinite sequence of real zeros when $m^2 < 0$; and that it has at most one real zero for $m^2 > 0$. From these results it follows that when $J > \frac{1}{4}$ and $m^2 < 0$, there exists, for each wave number $\kappa$, an infinite set of distinct, real characteristic values for $\nu$; and, when $J < \frac{1}{4}$ and $m^2 > 0$, there exists, again, for each wave number $\kappa$, either one real characteristic value for $\nu$ or none. Thus, when $J > \frac{1}{4}$, the flow would appear to be stable; and when $J < \frac{1}{4}$, while the flow does not appear unstable, the solution of the characteristic value problem, as formulated, does not seem to lead to a complete set of proper values and proper functions. This paradoxical situation was clarified by Case who pointed out that the solutions of equation~\eqref{eq:11-86} do not, by any means, exhaust the possible solutions: for, the equation
\begin{equation}
\label{eq:11-92}
x^2 f(x) = 0
\end{equation}
has in addition to the solution $f(x) = 0$, the solutions
\begin{equation}
\label{eq:11-93}
f(x) = \delta(x) \quad\text{and}\quad f(x) = \delta'(x).
\end{equation}
Therefore, we must supplement the solution~\eqref{eq:11-87} by the solutions of the equations
\begin{equation}
\label{eq:11-94}
\mathscr{L}W = \delta(\zeta) \quad\text{and}\quad \mathscr{L}W = \delta'(\zeta).
\end{equation}

The solutions of these latter equations can be expressed in terms of suitably defined Green's functions. It is found that the corresponding characteristic values form real continua. The existence of these continua of characteristic values means that in the expansion of an arbitrary perturbation into normal modes, we must include the proper functions derived as solutions of equations~\eqref{eq:11-94} along with those of equation~\eqref{eq:11-86}. And when these are included, the result, as Case has shown, is the following: for no positive $J$ can an arbitrary perturbation grow indefinitely with time; for $J > \frac{1}{4}$, an initial perturbation becomes a sum of oscillatory terms plus a term (derived from the continua) which tends to zero as $t^{\frac{1}{2}}$; for $J < \frac{1}{4}$, the corresponding asymptotic behaviour is $t^{-\frac{3}{2}+\sqrt{1/4-J}}$. It would, therefore, be strictly correct to say that \textit{for the distributions~\eqref{eq:11-75} extending over a semi-infinite region, the flow is stable for all Richardson numbers}.

\section{An example of the instability of a shear layer in an unbounded heterogeneous inviscid fluid}\label{sec:104}

The analysis of Dyson and Case, summarized in \S\,103\,(a), has shown that a linear variation of the stream velocity in a semi-infinite heterogeneous fluid with an exponentially decreasing density is, strictly, stable for all values of the Richardson number. While one cannot, of course, be certain of this, one is left with the suspicion that the result may in part derive from the assumed linear variation of the stream velocity $U$ which, on the one hand, eliminates a principal term from equation~\eqref{eq:11-22} (namely, the one in $D^2U$) and, on the other, makes $U$ unbounded. On these grounds, one may wonder if the result would not have been different had $D^2U$ been different from zero and $U$ been bounded.

A remarkable example to which the suspicions we have referred do not apply, and which at the same time allows an elementary and an explicit solution, has recently been discovered by Drazin. Drazin's example is
\begin{equation}
\label{eq:11-95}
U = U_0\tanh(z/d),
\end{equation}
where $U_0$ and $d$ are constants. On this distribution
\begin{equation}
\label{eq:11-96}
U \to +U_0 \;\text{as}\; z \to +\infty \quad\text{and}\quad U \to -U_0 \;\text{as}\; z \to -\infty.
\end{equation}

Setting $k_x = k$ in equation~\eqref{eq:11-22} and ignoring the effect of the heterogeneity on the inertia (by omitting the last term in $D\rho/\rho$), we have
\begin{equation}
\label{eq:11-97}
(n/k+U)(D^2-k^2)w-(D^2U)w-\frac{D\rho}{\rho}\frac{gk^2}{n/k+U}\frac{w}{n/k+U} = 0.
\end{equation}

We shall again suppose that
\begin{equation}
\label{eq:11-98}
\rho = \rho_0\,e^{-\beta z} \quad\text{and}\quad D\rho/\rho = -\beta = \text{constant}.
\end{equation}

For this distribution of the density, equation~\eqref{eq:11-97} takes the form
\begin{equation}
\label{eq:11-99}
(U-c)(D^2-k^2)w-(D^2U)w+J\frac{w}{U-c} = 0,
\end{equation}
if we measure length and velocity in the units $d$ and $U_0$, respectively, and put
\begin{equation}
\label{eq:11-100}
c = -n/kU_0 \quad\text{and}\quad J = g\beta d^2/U_0^2.
\end{equation}

In the chosen units,
\begin{equation}
\label{eq:11-101}
U = \tanh z;
\end{equation}
and we shall seek solutions of equation~\eqref{eq:11-99} which satisfy the boundary conditions
\begin{equation}
\label{eq:11-102}
w \to 0 \quad\text{as } z \to 0 \quad\text{and}\quad z \to \pm\infty.
\end{equation}
(The infinity in $\rho$, as $z \to -\infty$, would not seem to be relevant since the phenomenon we are concerned with is essentially limited to the neighbourhood $|z| \leqslant 1$.)

Returning to equation~\eqref{eq:11-99}, we may first observe that if $c$ is a characteristic value belonging to a proper solution $w$ and assigned $k$ and $J$, then $c^*$ is also a characteristic value for the same $k$ and $J$ and belongs to $w^*$. Moreover, if $U$ is an odd function of $z$ (as is the case with~\eqref{eq:11-101}) then $-c$ is also a characteristic value and belongs to $w(-z)$. On these symmetry grounds, it would appear that the principle of the exchange of stabilities is valid and that the state of marginal stability is characterized by $c = 0$. Assuming, then, that this is the case, we consider the equation
\begin{equation}
\label{eq:11-103}
U(D^2-k^2)w-(D^2U)w+J\frac{w}{U} = 0.
\end{equation}

We shall first transform equation~\eqref{eq:11-103} by using $U$ as the independent variable. Making use of the relations
\begin{equation}
\label{eq:11-104}
\frac{dU}{dz} = \operatorname{sech}^2 z = 1-U^2 \quad\text{and}\quad \frac{d^2U}{dz^2} = -2U(1-U^2),
\end{equation}
we find that equation~\eqref{eq:11-103} becomes
\begin{equation}
\label{eq:11-105}
\frac{d}{dU}\!\left[(1-U^2)\frac{dw}{dU}\right]+\left[2-\frac{k^2}{1-U^2}+\frac{J}{U^2(1-U^2)}\right]w = 0;
\end{equation}
and the corresponding boundary conditions are
\begin{equation}
\label{eq:11-106}
w = 0 \quad\text{for } U = \pm 1.
\end{equation}

The singular points of equation~\eqref{eq:11-105} are
\begin{equation}
\label{eq:11-107}
U = \pm 1 \quad\text{and}\quad U = 0.
\end{equation}

The behaviour of the solutions of equation~\eqref{eq:11-105} at these singular points is readily determined. Thus, by putting $U = 1+y$ and treating $y$ as small, we obtain the equation
\begin{equation}
\label{eq:11-108}
y\frac{d^2w}{dy^2}+\frac{dw}{dy}+\frac{1}{4y}(k^2-J)w = 0.
\end{equation}

This equation allows a solution of the form
\begin{equation}
\label{eq:11-109}
w = y^{\nu},
\end{equation}
where
\begin{equation}
\label{eq:11-110}
\nu^2 = \tfrac{1}{4}(k^2-J).
\end{equation}
Therefore, a solution of equation~\eqref{eq:11-105}, regular at $U = +1$ and $-1$, must have the behaviour
\begin{equation}
\label{eq:11-111}
w \to \text{constant}(1-U^2)^{\nu} \quad (U \to \pm 1),
\end{equation}
where
\begin{equation}
\label{eq:11-112}
\nu = \tfrac{1}{2}\sqrt{k^2-J};
\end{equation}
and it is necessary that the real part of $\nu$ is non-negative. In a similar manner, we find that the indicial equation, which determines the behaviour of $w$ at $U = 0$, is
\begin{equation}
\label{eq:11-113}
\mu^2-\mu+J = 0.
\end{equation}
Therefore, a solution of equation~\eqref{eq:11-105}, regular at $U = 0$, must have the behaviour
\begin{equation}
\label{eq:11-114}
w \to \text{constant}\,U^{\mu},
\end{equation}
where
\begin{equation}
\label{eq:11-115}
\mu = \tfrac{1}{2}+\tfrac{1}{2}\sqrt{1-4J}.
\end{equation}

The required behaviours~\eqref{eq:11-111} and~\eqref{eq:11-114} of $w$ suggest that we seek a solution of the form
\begin{equation}
\label{eq:11-116}
w = U^{\mu}(1-U^2)^{\nu}\chi,
\end{equation}
where $\chi$ is regular at $U = 0$, $+1$, and $-1$. For the more general substitution
\begin{equation}
\label{eq:11-117}
w = f\chi
\end{equation}
equation~\eqref{eq:11-105} becomes
\begin{multline}
\label{eq:11-118}
\frac{d^2\chi}{dU^2}+\left[\frac{1}{f}\frac{df}{dU}-\frac{2U}{1-U^2}\right]\frac{d\chi}{dU} \\
+\frac{1}{f}\!\left[\frac{d^2f}{dU^2}-\frac{2U}{1-U^2}\frac{1}{f}\frac{df}{dU}-\frac{k^2}{U^2(1-U^2)}+\frac{2U}{1-U^2}\frac{1}{f}\frac{df}{dU}\right]\chi = 0.
\end{multline}

For the particular substitution~\eqref{eq:11-116},
\begin{equation}
\label{eq:11-119}
f = U^{\mu}(1-U^2)^{\nu},
\end{equation}
and we readily verify that
\begin{equation}
\label{eq:11-120}
\frac{1}{f}\frac{df}{dU} = \frac{\mu}{U}-\frac{2\nu U}{1-U^2}
\end{equation}
and
\begin{equation}
\label{eq:11-121}
\frac{1}{f}\frac{df}{dU} = \frac{\mu(\mu-1)}{U^2(1-U^2)^{\nu}}+\frac{U^2}{(1-U^2)^2}\!\left[\frac{2\nu(2\nu+1)}{1-U^2}+\mu(\mu-1)\right]\frac{1}{U^2}.
\end{equation}

Inserting these expressions in equation~\eqref{eq:11-118} and making use of the relations~\eqref{eq:11-110} and~\eqref{eq:11-113} in the reductions, we find that the equation takes the simple form
\begin{equation}
\label{eq:11-122}
\frac{d^2\chi}{dU^2}+\left[\frac{\mu}{U}-\frac{2(\nu+1)U}{1-U^2}\right]\frac{d\chi}{dU}-\frac{(2\nu+\mu+2)(2\nu+\mu-1)}{1-U^2}\chi = 0.
\end{equation}

It is a remarkable fact noticed by Drazin that the special solution,
\begin{equation}
\label{eq:11-123}
\chi = \text{constant},
\end{equation}
and
\begin{equation}
\label{eq:11-124}
(2\nu+\mu+2)(2\nu+\mu-1) = 0,
\end{equation}
which we can write down by simply looking at equation~\eqref{eq:11-122}, provides the solution for the problem on hand.

Of the two factors in~\eqref{eq:11-124}, the first of them
\begin{equation}
\label{eq:11-125}
2\nu+\mu+2 = \sqrt{\kappa^2-J}+\tfrac{1}{2}\sqrt{1-4J}+2 \neq 0;
\end{equation}
hence, it is the second factor which must vanish. Thus,
\begin{equation}
\label{eq:11-126}
2\nu+\mu-1 = \sqrt{k^2-J}-\tfrac{1}{2}+\tfrac{1}{2}\sqrt{1-4J} = 0,
\end{equation}
or
\begin{equation}
\label{eq:11-127}
\sqrt{k^2-J} = \tfrac{1}{2}\!\left[1-\sqrt{1-4J}\right]\!.
\end{equation}
On squaring relation~\eqref{eq:11-127}, we obtain
\begin{equation}
\label{eq:11-128}
\sqrt{1-4J} = 1-2k^2;
\end{equation}
and on squaring this relation, we obtain
\begin{equation}
\label{eq:11-129}
J = k^2(1-k^2).
\end{equation}

The corresponding expressions for $\nu$ and $\mu$ are
\begin{equation}
\label{eq:11-130}
\nu = \tfrac{1}{2}k^2 \quad\text{and}\quad \mu = 1-k^2,
\end{equation}
so that the proper function which belongs to the characteristic value~\eqref{eq:11-129} is
\begin{equation}
\label{eq:11-131}
w = \text{constant } (1-U^2)^{k^2/2}U^{1-k^2};
\end{equation}
or, reverting to the variable $z$, we have
\begin{equation}
\label{eq:11-132}
w = \text{constant}\,(\operatorname{sech} z)^{k^2}(\tanh z)^{1-k^2}.
\end{equation}

[Notice that when $J = 0$ equation~\eqref{eq:11-105} becomes
\begin{equation}
\label{eq:11-133}
\frac{d}{dU}\!\left[(1-U^2)\frac{dw}{dU}\right]+\left[2-\frac{k^2}{1-U^2}\right]w = 0.
\end{equation}
The solutions of this equation are the associated Legendre functions of degree 1 and order $k$. Only for $k^2 = 0$ and $1$ does equation~\eqref{eq:11-133} allow solutions which satisfy the boundary conditions~\eqref{eq:11-106}; and, in these special cases, the solutions are
\begin{equation}
\label{eq:11-134}
w = \text{constant}\,U \quad (k^2 = 0)
\end{equation}
and
\begin{equation}
w = \text{constant}\,\sqrt{1-U^2} \quad (k^2 = 1).
\end{equation}
It will be observed that these solutions are included in~\eqref{eq:11-131}.]

\figurePlaceholder{120}{The regions of stability and instability in the $(J,k)$ plane for a shear layer with the velocity distribution $U = U_0\tanh(z/d)$; $k$ is the wave number measured in the unit $1/d$ and $J$ is the Richardson number.}

The curve of marginal stability defined by equation~\eqref{eq:11-129} is shown in Fig.\ 120. The possibility of closed unstable regions above the curve, and similar stable regions below the curve, cannot be excluded; but it does not seem likely on energetic considerations such as those which led to the condition~\eqref{eq:11-74}. It would appear, then, that the Richardson number defining the limit of instability in the $(J,k)$ plane is determined by equation~\eqref{eq:11-129}; it is given by
\begin{equation}
\label{eq:11-135}
J_{\max} = \tfrac{1}{4} \quad\text{for } k^2 = \tfrac{1}{2}.
\end{equation}
This result is in exact agreement with~\eqref{eq:11-74} and lends support to the physical arguments which were advanced in \S\,103.

\section{The effect of rotation}\label{sec:105}

We shall now consider the effect of rotation on the development of the Kelvin--Helmholtz instability. The initial stationary state is, then, the same as that postulated in \S\,100, except that the fluid will now be supposed to be in uniform rotation about the $z$-axis with an angular velocity $\Omega$. Under these circumstances, equations~\eqref{eq:11-1} and~\eqref{eq:11-2} will be replaced by
\begin{equation}
\label{eq:11-136}
\rho\frac{\partial u}{\partial t}+\rho U\frac{\partial u}{\partial x}+\rho w\frac{dU}{dz}-2\rho\Omega v = -\frac{\partial}{\partial x}\,\delta p,
\end{equation}
\begin{equation}
\label{eq:11-137}
\rho\frac{\partial v}{\partial t}+\rho U\frac{\partial v}{\partial x}+2\rho\Omega u = -\frac{\partial}{\partial y}\,\delta p,
\end{equation}
while equations~\eqref{eq:11-3}--\eqref{eq:11-6} remain unaltered. For solutions having the $(x,y,t)$-dependence given by~\eqref{eq:11-7}, equations~\eqref{eq:11-136} and~\eqref{eq:11-137} give
\begin{equation}
\label{eq:11-138}
\mathrm{i}\rho(n+k_x U)u+\rho(DU)w-2\rho\Omega v = -\mathrm{i}k_x\,\delta p
\end{equation}
and
\begin{equation}
\label{eq:11-139}
\mathrm{i}\rho(n+k_x U)v+2\rho\Omega u = -\mathrm{i}k_y\,\delta p.
\end{equation}
From equations~\eqref{eq:11-138} and~\eqref{eq:11-139} we obtain (in place of equation~\eqref{eq:11-14})
\begin{equation}
\label{eq:11-140}
\mathrm{i}\rho(n+k_x U)Dw - \mathrm{i}\rho k_x(DU)w + 2\rho\Omega\zeta = -k^2\delta p,
\end{equation}
where
\begin{equation}
\label{eq:11-141}
\zeta = \mathrm{i}k_x v - \mathrm{i}k_y u
\end{equation}
is the $z$-component of the vorticity. By multiplying equations~\eqref{eq:11-138} and~\eqref{eq:11-139} by $-\mathrm{i}k_y$ and $+\mathrm{i}k_x$, respectively, adding, and making use of the equation of continuity~\eqref{eq:11-13}, we also obtain
\begin{equation}
\label{eq:11-142}
\mathrm{i}\rho(n+k_x U)\zeta = \mathrm{i}k_x\rho(DU)w + 2\rho\Omega Dw.
\end{equation}

Substituting for $\zeta$ from this equation in equation~\eqref{eq:11-140}, we obtain after some rearrangement
\begin{equation}
\label{eq:11-143}
\mathrm{i}\rho(n+k_x U)\!\left[1-\frac{4\Omega^2}{(n+k_x U)^2}\right]\!Dw - \mathrm{i}\rho\!\left(k_x+\frac{2\mathrm{i}k_x\Omega}{n+k_x U}\right)\!(DU)w = -k^2\delta p.
\end{equation}

Now eliminating $\delta p$ from equation~\eqref{eq:11-15} (which continues to be valid) and equation~\eqref{eq:11-143}, we obtain
\begin{multline}
\label{eq:11-144}
D\!\left[\rho(n+k_x U)\!\left(1-\frac{4\Omega^2}{(n+k_x U)^2}\right)\!(DU)w\right] \\
-k^2\rho(n+k_x U)w = gk^2\bigg[D\rho - \frac{k^2}{g}\sum_s T_s\,\delta(z-z_s)\bigg]\frac{w}{n+k_x U}.
\end{multline}

As in \S\,100, we must require that at the interface between two fluids $w/(n+k_x U)$ is continuous. By integrating equation~\eqref{eq:11-144} across such an interface, we obtain the further condition
\begin{multline}
\label{eq:11-145}
\Delta_s\!\left[\rho(n+k_x U)\!\left(1-\frac{4\Omega^2}{(n+k_x U)^2}\right)\!(DU)w\right] \\
- \rho\!\left(k_x+\frac{2\mathrm{i}k_x\Omega}{n+k_x U}\right)\!\Delta_s\!\left(\rho\right)\!\left(\frac{w}{n+k_x U}\right)_{s} \\
= gk^2\!\left[\Delta_s(\rho)-\frac{k^2 T_s}{g}\right]\!\bigg|_{}\frac{w}{n+k_x U}\bigg|_s\!,
\end{multline}
where $\Delta_s(f)$ denotes the jump which a quantity $f$ experiences at the interface $z = z_s$; and the subscript $a$ signifies the value which a quantity, known to be continuous at an interface, takes at $z = z_s$.

\subsection*{(a) The case of two uniform fluids in relative horizontal motion separated by a horizontal boundary}

Consider the case of two uniform fluids of densities $\rho_1$ and $\rho_2$ separated by a horizontal boundary at $z = 0$; and let the velocities of streaming in the two fluids be $U_1$ and $U_2$. In the two regions of constant $\rho$ and $U$, equation~\eqref{eq:11-144} becomes
\begin{equation}
\label{eq:11-146}
\left[1-\frac{4\Omega^2}{(n+k_x U)^2}\right]D^2w = k^2 w,
\end{equation}
where, for the present, we are suppressing the subscripts distinguishing the two fluids. The general solution of equation~\eqref{eq:11-146} is a linear combination of the integrals
\begin{equation}
\label{eq:11-147}
e^{+\kappa z} \quad\text{and}\quad e^{-\kappa z},
\end{equation}
where
\begin{equation}
\label{eq:11-148}
\kappa = \frac{k}{\sqrt{1-4\Omega^2/(n+k_x U)^2}}.
\end{equation}

Since $w$ must vanish when $z \to \pm\infty$ and $w/(n+k_x U)$ must be continuous at $z = 0$, we may write
\begin{equation}
\label{eq:11-149}
w_1 = A(n+k_x U_1)e^{+\kappa_1 z} \quad (z < 0)
\end{equation}
and
\begin{equation}
\label{eq:11-150}
w_2 = A(n+k_x U_2)e^{-\kappa_2 z} \quad (z > 0)
\end{equation}
as the solutions appropriate in the two regions, provided, in the definitions of $\kappa_1$ and $\kappa_2$, the signs of the square roots are so chosen that the real parts of $\kappa_1$ and $\kappa_2$ are positive.

When $DU = 0$ on either side of the interface (as in the present instance), the condition~\eqref{eq:11-145}, which must obtain generally reduces to
\begin{equation}
\label{eq:11-151}
\Delta_s\!\left[\rho(n+k_x U)\frac{Dw}{\kappa}\right] = g\!\left[\Delta_s(\rho)-\frac{k^2T}{g}\right]\!\frac{w}{(n+k_x U)_s}.
\end{equation}

Applying this condition to the solutions given by~\eqref{eq:11-149} and~\eqref{eq:11-150}, we find
\begin{equation}
\label{eq:11-152}
\frac{\rho_1}{\kappa_1}(n+k_x U_1)^2+\frac{\rho_2}{\kappa_2}(n+k_x U_2)^2 = g\!\left(\rho_1-\rho_2+\frac{k^2T}{g}\right)\!.
\end{equation}

Substituting for $\kappa_1$ and $\kappa_2$ in accordance with~\eqref{eq:11-148} and letting $\alpha_1$ and $\alpha_2$ have the meanings assigned in~\eqref{eq:11-27}, we obtain the characteristic equation
\begin{multline}
\label{eq:11-153}
\alpha_1(n+k_x U_1)^2\!\left[1-\frac{4\Omega^2}{(n+k_x U_1)^2}\right]^{1/2}+\alpha_2(n+k_x U_2)^2\!\left[1-\frac{4\Omega^2}{(n+k_x U_2)^2}\right]^{1/2} \\
= gk\!\left[(\alpha_1-\alpha_2)+\frac{k^2T}{g(\rho_1+\rho_2)}\right]\!.
\end{multline}

When $k_x = 0$, equation~\eqref{eq:11-153} reduces to equation X~(160) appropriate for Rayleigh--Taylor instability in the presence of rotation.\footnote{It may be recalled here that what was denoted by $n$ in Chapter~X is now denoted by $w$.} Clearly, the Kelvin--Helmholtz instability is least uninhibited for perturbations in the direction of streaming. Accordingly, replacing $k_x$ by $k$ in equation~\eqref{eq:11-153}, we have
\begin{multline}
\label{eq:11-154}
\alpha_1(n+k U_1)^2\!\left[1-\frac{4\Omega^2}{(n+k U_1)^2}\right]^{1/2}+\alpha_2(n+k U_2)^2\!\left[1-\frac{4\Omega^2}{(n+k U_2)^2}\right]^{1/2} \\
= gk\!\left[(\alpha_1-\alpha_2)+\frac{k^2T}{g(\rho_1+\rho_2)}\right]\!.
\end{multline}

\subsection*{(b) The discussion of the characteristic equation}

We shall restrict our discussion of equation~\eqref{eq:11-154} to the case $T = 0$; and consider, first, the stable modes of oscillation which are possible by a graphical method used by Taylor in a different connexion.

Ignoring the term in the surface tension, we can rewrite equation~\eqref{eq:11-154} in the form
\begin{equation}
\label{eq:11-155}
\alpha_1(U_1-c)^2\!\left[1-\frac{4\Omega^2}{k^2(U_1-c)^2}\right]^{1/2}+\alpha_2(U_2-c)^2\!\left[1-\frac{4\Omega^2}{k^2(U_2-c)^2}\right]^{1/2} = \frac{g}{k}(\alpha_1-\alpha_2),
\end{equation}
where we have put
\begin{equation}
\label{eq:11-156}
c = -n/k.
\end{equation}

By measuring velocities in the unit $\sqrt{g/k}$ and letting
\begin{equation}
\label{eq:11-157}
U_1-c = \xi\sqrt{g/k} \quad\text{and}\quad U_2-c = \eta\sqrt{g/k},
\end{equation}
we have
\begin{equation}
\label{eq:11-158}
\frac{\alpha_1}{\alpha_1-\alpha_2}\,\xi^2\sqrt{1-\frac{\omega^2}{\xi^2}}+\frac{\alpha_2}{\alpha_1-\alpha_2}\,\eta^2\sqrt{1-\frac{\omega^2}{\eta^2}} = 1,
\end{equation}
where
\begin{equation}
\label{eq:11-159}
\omega^2 = 4\Omega^2/gk.
\end{equation}

It may be recalled here that restriction on the sign of the square root in the definition of $\kappa$ in equation~\eqref{eq:11-148} requires that in equation~\eqref{eq:11-158} the signs of the square roots must be so chosen that their real parts are positive.

So long as $\xi$ and $\eta$ are real and $|\xi|$ and $|\eta|$ are each greater than $\omega$, equation~\eqref{eq:11-158} defines a real locus in the $(\xi,\eta)$-plane. This locus consists of four disjointed curves (the parts designated by $P$ in Fig.~121\,$a,b$) symmetrically disposed with respect to the axes. The part in the first quadrant connects the points
\begin{equation}
\label{eq:11-160}
(\omega,\eta_0) \quad\text{and}\quad (\xi_0,\omega),
\end{equation}
where $\xi_0$ and $\eta_0$ are determined by the equations
\begin{equation}
\label{eq:11-161}
\alpha_1\,\xi_0^2\sqrt{1-\frac{\omega^2}{\xi_0^2}} = \alpha_1-\alpha_2 \quad\text{and}\quad \alpha_2\,\eta_0^2\sqrt{1-\frac{\omega^2}{\eta_0^2}} = \alpha_1-\alpha_2.
\end{equation}

The explicit expressions for $\xi_0$ and $\eta_0$ are
\begin{equation}
\label{eq:11-162}
\xi_0 = \left\{\tfrac{1}{2}\omega^2+\tfrac{1}{2}\!\left[\omega^4+4\!\left(\frac{\alpha_1-\alpha_2}{\alpha_1}\right)^{\!2}\right]^{1/2}\right\}^{1/2}
\end{equation}
and
\begin{equation}
\label{eq:11-163}
\eta_0 = \left\{\tfrac{1}{2}\omega^2+\tfrac{1}{2}\!\left[\omega^4+4\!\left(\frac{\alpha_1-\alpha_2}{\alpha_2}\right)^{\!2}\right]^{1/2}\right\}^{1/2}\!.
\end{equation}
Since $\alpha_1 > \alpha_2$,
\[
\eta_0 > \xi_0 > \omega.
\]

\figurePlaceholder{121}{The loci defined by the equations $\alpha_1\xi^2\sqrt{1-\omega^2/\xi^2}\pm\alpha_2\eta^2\sqrt{1-\omega^2/\eta^2} = \alpha_1-\alpha_2$. When the signs are dissimilar they define the dashed-line curves; the branch designated $S_2$ corresponds to the first sign in the equation being positive and the second sign negative; and the branch designated $S_1$ corresponds to the first sign in the equation being negative and the second sign positive. The branches $S_1$ and $S_2$ are connected by the branch $P$ which corresponds to both signs in the equation being positive. The cases illustrated are for $\alpha_1=\alpha_2 = 0{\cdot}5$ and $\omega^2 = 0{\cdot}5$ (the figure on the right) and $\omega^2 = 0{\cdot}1$ (the figure on the left).}

It is also clear that at $(\omega,\eta_0)$ the locus has a vertical tangent, while at $(\xi_0,\omega)$ it has a horizontal tangent.

In Table~LII, $\xi_0$ and $\eta_0$ are listed for a few values of $\omega$ for the case $\alpha_1-\alpha_2 = 0{\cdot}5$.

\begin{table}[ht]
\centering
\caption{The coordinates of the end points $\xi_0$, $\omega$ and $\omega$, $\eta_0$ of the $(\xi,\eta)$-locus ($\alpha_1-\alpha_2 = 0{\cdot}5$)}
\label{tab:LII}
\smallskip
\begin{tabular}{c|cc|cc|cc|cc}
\hline
& \multicolumn{2}{c|}{$l=2$} & \multicolumn{2}{c|}{$l=3$} & \multicolumn{2}{c|}{$l=4$} \\
$\omega^2$ & $\xi_0$ & $\eta_0$ & $\xi_0$ & $\eta_0$ & $\xi_0$ & $\eta_0$ \\
\hline
0 & $\alpha_1$ & $\alpha_2$ & & & & \\
0.1 & 0.8164 & 1.4142 & 0.8 & 1 & 1.1568 & 5.0005 \\
0.2 & 0.8196 & 1.4150 & 0.8 & 1.0 & 1.1568 & 5.0005 \\
0.3 & 0.8399 & 1.4213 & & & 1.2945 & 1.6569 \\
0.4 & 0.8476 & 1.4330 & 1.0 & 1.2 & 1.4715 & 1.7999 \\
0.5 & 0.8656 & 1.4498 & 1.0 & 1.2 & 1.5601 & 1.9177 \\
0.6 & 0.9331 & 1.4792 & 1.3 & 1.5 & 1.8351 & 2.0429 \\
0.7 & 0.9868 & 1.5054 & & & & \\
0.8 & 1.0547 & 1.5866 & & & & \\
0.9 & 1.2350 & 1.7639 & & & & \\
1.0 & 1.4359 & 1.9552 & 1.9 & 2.0 & 2.9169 & 2.9740 \\
\hline
\end{tabular}
\end{table}

For a given difference $V$ in the streaming velocities $U_1$ and $U_2$, $\xi$ and $\eta$ must lie on the line
\begin{equation}
\label{eq:11-164}
\xi-\eta = V\sqrt{k/g}.
\end{equation}

In order to find whether stable modes are possible for a given $V\sqrt{k/g}$, we have only to determine whether the corresponding line~\eqref{eq:11-164} has any intersections with the real part of the $(\xi,\eta)$-locus. Similarly, the effect of varying $V$ (keeping $k$ fixed) can be found by drawing a series of parallel lines inclined to the $\xi$-axis by $45^\circ$ and determining when intersections with the $(\xi,\eta)$-locus occur.

While the allowed real modes can be enumerated and studied in the manner we have described, it is more important to be able to determine when equations~\eqref{eq:11-158} and~\eqref{eq:11-164} permit a complex root and entail instability. In this latter connexion, particular interest attaches to the maximum value of $|V|$ (for a given $k$) such that for all $|V|$ less than this maximum, we can be assured of stability.

The problem of ascertaining the number and the character of the roots of equations~\eqref{eq:11-158} and~\eqref{eq:11-164} is not a straightforward one: the difficulty arises from the curious restriction on the signs of the square roots in equation~\eqref{eq:11-158}, namely, that the real parts must be positive. This latter restriction makes the problem non-algebraic and precludes the application of the fundamental theorem of algebra. However, we can restore to the problem its algebraic character by considering together with equation~\eqref{eq:11-164} the following set of four equations:
\begin{align}
\label{eq:11-165}
+\alpha_1\,\xi^2\!\sqrt{1-\omega^2/\xi^2}+\alpha_2\,\eta^2\!\sqrt{1-\omega^2/\eta^2} &= \alpha_1-\alpha_2, \\
\label{eq:11-166}
+\alpha_1\,\xi^2\!\sqrt{1-\omega^2/\xi^2}-\alpha_2\,\eta^2\!\sqrt{1-\omega^2/\eta^2} &= \alpha_1-\alpha_2, \\
\label{eq:11-167}
-\alpha_1\,\xi^2\!\sqrt{1-\omega^2/\xi^2}+\alpha_2\,\eta^2\!\sqrt{1-\omega^2/\eta^2} &= \alpha_1-\alpha_2, \\
\label{eq:11-168}
-\alpha_1\,\xi^2\!\sqrt{1-\omega^2/\xi^2}-\alpha_2\,\eta^2\!\sqrt{1-\omega^2/\eta^2} &= \alpha_1-\alpha_2.
\end{align}
(with the same restriction, as before, on the choice of the signs of the square roots). The four equations (together with~\eqref{eq:11-164}) constitute a pure algebraic problem; for, they are together equivalent to the single eighth degree polynomial equation which can be derived from any one of the four equations~\eqref{eq:11-165}--\eqref{eq:11-168} by clearing the square roots by two successive squarings. This equivalence is a complete and a reciprocal one: a root of any of the four equations is a root of the parent eighth degree equation; and any root of the parent equation must be a root of \textit{some} one of the four equations. The fundamental theorem of algebra, therefore, applies to the equations~\eqref{eq:11-164}--\eqref{eq:11-168} together; and we can conclude that \textit{they must have a total of exactly eight roots under all circumstances}. In addition to this result, there are a number of properties of the equations which we shall now enumerate. (1) If any of the four equations allows a complex root, it must also allow its complex conjugate as a root, i.e.\ each of the equations can allow complex roots only in conjugate pairs. (2) The roots which follow from any of the four equations~\eqref{eq:11-165}--\eqref{eq:11-168} are continuous functions of the parameters of the equation except at the singular points $(\pm\xi_0, \pm\omega)$ and $(\pm\omega, \pm\eta_0)$. (3) Only at the singular points can the number of roots, which any of them allows, change. (4) At a singular point only those equations can `exchange' roots which coalesce at that point. (This last remark will be clarified presently.)

We shall now show how the properties of the system~\eqref{eq:11-164}--\eqref{eq:11-168} which we have enumerated enable us to ascertain in a systematic way the number and the character of the roots of each of the four equations for any assigned value of $|V|$.

We have already seen how equation~\eqref{eq:11-165} defines in the $(\xi,\eta)$-plane a locus consisting of four disjointed pieces. We shall call this locus the principal or the $P$-branch of the equations~\eqref{eq:11-165}--\eqref{eq:11-168}. Similarly we shall call the loci defined by equations~\eqref{eq:11-166} and~\eqref{eq:11-167} the subsidiary branches $S_1$ and $S_2$, respectively. Equation~\eqref{eq:11-168} cannot be satisfied (with the imposed restriction on the signs of the square roots) by any real pair of values of $\xi$ and $\eta$; the equation accordingly defines a purely complex branch which we shall designate by $C$. On this last account no exchange of roots between this equation and the remaining three can take place. Consequently, the principal and the two subsidiary branches must conserve their total number of roots. We shall presently see that this number is six.

Now consider the equations~\eqref{eq:11-164}--\eqref{eq:11-168} for the special case
\begin{equation}
\label{eq:11-169}
V = 0 \quad\text{and}\quad \xi = \eta.
\end{equation}
The equations become
\begin{align}
\label{eq:11-170}
\xi^2\!\sqrt{1-\omega^2/\xi^2} &= \alpha_1-\alpha_2 && (P), \\
\label{eq:11-171}
\xi^2(1-\omega^2/\xi^2) &= 1 && (S_1), \\
\label{eq:11-172}
-\xi^2\!\sqrt{1-\omega^2/\xi^2} &= 1 && (S_2), \\
\label{eq:11-173}
-\xi^2(1-\omega^2/\xi^2) &= \alpha_1-\alpha_2 && (C).
\end{align}

By squaring these equations, we obtain
\begin{equation}
\label{eq:11-174}
\xi^4-\omega^2\xi^2-(\alpha_1-\alpha_2)^2 = 0 \quad (P \text{ and } C)
\end{equation}
and
\begin{equation}
\label{eq:11-175}
\xi^4-\omega^2\xi^2-1 = 0 \quad (S_1 \text{ and } S_2).
\end{equation}

The roots of these equations are
\begin{equation}
\label{eq:11-176}
\xi^2 = \tfrac{1}{2}\omega^2+\tfrac{1}{2}\sqrt{\omega^4+4(\alpha_1-\alpha_2)^2} \quad (P \text{ and } C)
\end{equation}
and
\begin{equation}
\label{eq:11-177}
\xi^2 = \tfrac{1}{2}\omega^2+\tfrac{1}{2}\sqrt{\omega^4+4} \quad (S_1 \text{ and } S_2).
\end{equation}

By substituting these roots back into the original equations~\eqref{eq:11-170}--\eqref{eq:11-173} we directly verify that the positive sign in equation~\eqref{eq:11-176} goes with equation~\eqref{eq:11-170} while the negative sign goes with equation~\eqref{eq:11-173}. Similarly, the positive sign in equation~\eqref{eq:11-177} goes with the $S_2$-branch while the negative sign goes with the $S_1$-branch. Therefore, for $V = 0$, the $P$- and the $S_1$-branches each allow two real roots, the $S_2$- and the $C$-branches each allow a conjugate pair of complex roots; this is in agreement with what we observe in Fig.\ 121.

Now let $|V| = |V|\sqrt{k/g}$ increase. The number and the character of the roots of the three branches, $P$, $S_1$, and $S_2$ must remain unchanged. At a singular point, such as $(\xi_0,\omega)$, only the branches $P$ and $S_1$ can (and do) exchange roots, while the branch $S_2$ is unaffected. In other words, at $(\xi_0,\omega)$ the total number of roots of the $P$- and $S_1$-branches must be conserved separately, while at $(\omega,\eta_0)$ the total number of roots of the $P$- and the $S_2$-branches must be conserved separately.

An important result which follows from the discussion up to this point is that the $P$-branch allows \textit{only} real roots (either one or two) so long as
\begin{equation}
\label{eq:11-178}
|V| \leqslant \eta_0-\omega;
\end{equation}
and for $|V|$ in excess of this value unstable modes are possible.

As $|V|$ increases beyond the limit set by~\eqref{eq:11-178}, the detailed sequence of events which follows depends on the precise value of $\omega$. For the case illustrated in Fig.~121\,(a), the sequence is as follows. We have seen that for $|V|$ just in excess of $(\eta_0-\omega)$ the $P$-branch has a complex pair of conjugate roots, the $S_2$-branch has one real root and a complex pair of roots, and the $S_1$-branch has one real root. The next change will occur when $|V| = (\xi_0+\omega)$. However, just prior to this, because of the horizontal tangency of the $(\xi,\eta)$-locus at $(\xi_0+\omega)$, the two complex roots of the $P$-branch briefly become a pair of real roots so that for a small range of $|V| < (\xi_0+\omega)$ there is a domain of stability. When $|V| = (\xi_0+\omega)$, because of the horizontal tangency of the $(\xi,\eta)$-locus at $(\pm\xi_0,\pm\omega)$, the $P$- and the $S_1$-branches coalesce and the number of roots of these two branches (equal to five prior to crossing) must be conserved. We observe that just after the crossing the $S_1$-branch has two real roots so that $P$ must have exchanged one of its two real roots (it had just prior to the crossing) with a complex pair. After this last crossing, therefore, the $P$-branch has one real root and the $S_1$-branch has two real roots, and the $S_2$-branch has one real root. The last change occurs at $|V| = (\eta_0+\omega)$ when the $P$- and the $S_2$-branches exchange roots. After the exchange, $S_2$ has two real roots and $P$ must be left with a pair of complex roots so that $P$ must have exchanged one of its two real roots.

\figurePlaceholder{122}{Figure 122.}

The sequence of events we have described beyond $|V| = (\eta_0-\omega)$ applies to the case when the loci have `opened out' in the manner of Fig.~121\,$a$. For smaller values of $\omega$ the situation is more complex. Thus for $\omega \to 0$ the number of cases which must be distinguished are shown in Fig.~122. But in principle there is no difficulty in continuing the enumeration we have described starting at $|V| = 0$.

The principal result established in this section is that \textit{for a system partaking in rotation, the Kelvin--Helmholtz instability cannot occur so long as}
\begin{equation}
\label{eq:11-179}
|V|\sqrt{\frac{k}{g}} \leqslant \left\{\tfrac{1}{2}\omega^2+\tfrac{1}{2}\!\left[\omega^4+4\!\left(\frac{\alpha_1-\alpha_2}{\alpha_2}\right)^{\!2}\right]^{1/2}\right\}^{1/2}\!\!-\omega.
\end{equation}

\section{The effect of a horizontal magnetic field}\label{sec:106}

% §106(a) content was blocked by content filter; equations (180)-(196) not available.
% §106(b)-(c) follows:

\subsection*{(b) The case of two uniform fluids in relative horizontal
  motion separated by a horizontal boundary}

We shall apply the foregoing equations to the case of two uniform
fluids of densities $\rho_1$ and $\rho_2$ separated by a horizontal boundary
at $z=0$.  Let the velocities of streaming of the two fluids be $U_1$
and~$U_2$.

In the two regions of constant $\rho$ and~$U$, the perturbation equations give
\begin{equation}
\label{eq:11-197}
\zeta = 0
\end{equation}
and
\begin{equation}
\label{eq:11-198}
\left[p(n+k_x\,U) - \frac{k_x^{2}}{n+k_x\,U}\,
\frac{\mu H^{2}}{4\pi}\right](D^{2}-k^{2})w = 0,
\end{equation}
where we have suppressed the subscripts distinguishing the two fluids.
The solutions in the two regions $z<0$ and $z>0$ can be written in the
forms (cf.\ equations~(24) and~(25))
\begin{equation}
\label{eq:11-199}
w_1 = A(n+k_x\,U_1)\,e^{+kz}\quad (z<0)
\end{equation}
and
\begin{equation}
\label{eq:11-200}
w_2 = A(n+k_x\,U_2)\,e^{-kz}\quad (z>0).
\end{equation}
By this choice of the solutions, the condition requiring the continuity
of $w/(n+k_x\,U)$, as well as the conditions at infinity, have been met.
It remains to satisfy the further condition which follows by integrating
across the interface at $z=0$.  This condition is, clearly,
\begin{equation}
\label{eq:11-201}
\Delta_0\!\bigl\{p(n+k_x\,U)Dw\bigr\}
= k_x^{2}\,\Delta_0\!\!\left(\frac{Dw}{n+k_x\,U}\right)
  + gk^{2}\Delta_0(p)\!\left(\frac{w}{n+k_x\,U}\right)_{\!0}.
\end{equation}

Now applying the condition~\eqref{eq:11-201} to the solutions given by
equations~\eqref{eq:11-199} and~\eqref{eq:11-200}, we obtain the
characteristic equation
\begin{equation}
\label{eq:11-202}
\rho_1(n+k_x\,U_1)^{2}+\rho_2(n+k_x\,U_2)^{2}
= gk(\rho_1-\rho_2)+k_x^{2}\,\frac{\mu H^{2}}{2\pi}.
\end{equation}

On comparing this equation with equation~(26), we observe that in this
instance the effect of the magnetic field is equivalent to that of a
surface tension of amount
\begin{equation}
\label{eq:11-203}
T_{\text{eff}} = \frac{\mu H^{2}}{2\pi}\,\frac{k_x^{2}}{k^{2}}.
\end{equation}

The roots of equation~\eqref{eq:11-202} are (cf.\ equation~(30))
\begin{equation}
\label{eq:11-204}
n = -k_x(\alpha_1\,U_1 + \alpha_2\,U_2) \pm
\left[gk(\alpha_1-\alpha_2) + k_x^{2}\,\frac{\mu H^{2}}{2\pi}
  - k_x^{2}\,\alpha_1\,\alpha_2(U_1-U_2)^{2}\right]^{1/2}\!.
\end{equation}

From this equation it follows that \textit{a uniform magnetic field acting
parallel to the direction of streaming will suppress the
Kelvin--Helmholtz instability if}
\begin{equation}
\label{eq:11-205}
\alpha_1\,\alpha_2\,(U_1-U_2)^{2}
\leqslant \frac{\mu H^{2}}{2\pi(\rho_1+\rho_2)},
\end{equation}
\textit{i.e.\ if the relative speed does not exceed the root-mean-square
Alfv\'en speed in the two media.}

\subsection*{(c) The effect of a magnetic field transverse to the
  direction of streaming}

When the impressed magnetic field is transverse to the direction
of~$U$, the relevant perturbation equations are
\begin{equation}
\label{eq:11-206}
\mathrm{i}p(n+k_x\,U)u + p(DU)w
+ \frac{\mu H}{4\pi}\,(\mathrm{i}k_x\,h_x - \mathrm{i}k_y\,h_y)
= -\mathrm{i}k_x\,\delta p,
\end{equation}
\begin{equation}
\label{eq:11-207}
\mathrm{i}p(n+k_x\,U)v = -\mathrm{i}k_y\,\delta p,
\end{equation}
\begin{equation}
\label{eq:11-208}
\mathrm{i}p(n+k_x\,U)w
- \frac{\mu H}{4\pi}\,(\mathrm{i}k_x\,h_x - Dh_x)
= -D\delta p - g\,\delta\rho,
\end{equation}
and
\begin{equation}
\label{eq:11-209}
\mathrm{i}(n+k_x\,U)\mathbf{h}
= \mathrm{i}k_x\,H\,\mathbf{u} + h_x\,DU.
\end{equation}

Treating the present set of equations in the same manner as in
\S\,(a) above, we now obtain
\begin{align}
\label{eq:11-210}
D\!\bigl\{p(n+k_x\,U)Dw - pk_x(DU)w\bigr\}
&= k^{2}p(n+k_x\,U)w \notag\\
&\quad + gk^{2}\,\frac{Dp}{n+k_x\,U}\,w.
\end{align}

We observe that equation~\eqref{eq:11-210} reduces to the hydrodynamic
equation~(16) of \S\,100 in case $k_x=0$.  Therefore, \textit{the
development of the Kelvin--Helmholtz instability, in the direction of
the streaming, is uninfluenced by the presence of the magnetic field in
the transverse direction.}

\section*{Bibliographical Notes}

The following general references may be noted:

\begin{enumerate}

\item \textsc{Sir Horace Lamb}, \emph{Hydrodynamics},
  \S\S\,232 and 268, Cambridge, England, 1932.

\item \textsc{J. Proudman}, \emph{Dynamical Oceanography},
  chap.\ xii, \S\,181--3, John Wiley \& Sons, Inc., New York, 1953.

\medskip
\noindent\S\,101. The first reference to what has come to be known as the
Kelvin--Helmholtz instability occurs in the writings of Helmholtz:

\item \textsc{H. Helmholtz}, `Ueber discontinuirliche Fl\"ussigkeitsbewegungen',
  \emph{Wissenschaftliche Abhandlungen}, 146--57, J.~A. Barth, Leipzig, 1882;
  translation by Guthrie in \emph{Phil.\ Mag.}, Ser.~4, \textbf{36}, 337--46 (1868).

\smallskip\noindent
(The quotation in the text is from Guthrie's translation.)

\smallskip\noindent
Helmholtz's discussion was largely qualitative. The same basic question was
treated analytically by:

\item \textsc{Lord Kelvin}, `Hydrokinetic solutions and observations',
  `On the motion of free solids through a liquid, 69--75',
  `Influence of wind and capillarity on waves in water supposed frictionless, 76--85',
  \emph{Mathematical and Physical Papers}, iv,
  \emph{Hydrodynamics and General Dynamics}, Cambridge, England, 1910.

\smallskip\noindent
Kelvin's discussion of the stability of superposed fluids in a state of
differential streaming is exceptionally complete; there is hardly a formula
in this section which cannot be found in Kelvin's paper. The critical wind
velocity, 650~cm/sec (or 660~cm/sec as Kelvin gave it), occurs, for example,
for the first time in this paper. As stated in the text, the significance of
this particular velocity for the generation of waves by wind has been the
subject of much discussion; an account of it can be found in Proudman's book
(reference~2; chap.~xii). The recognition that Kelvin's critical velocity
may, indeed, have a bearing on phenomena which are observed is more recent:

\item \textsc{W. H. Munk}, `A critical wind speed for air--sea boundary processes',
  \emph{J.\ Marine Research}, \textbf{6}, 203--18 (1947).


\smallskip\noindent
For a further discussion of these and related matters, see:

\item \textsc{J. W. Miles}, `On the generation of surface waves by shear flows',
  \emph{J.\ Fluid Mech.}, \textbf{3}, 185--204 (1957).

\item \makebox[2em][l]{------} `On the generation of surface waves by shear flows.~II',
  ibid.\ \textbf{6}, 568--82 (1959).

\item \makebox[2em][l]{------} `On the generation of surface waves by shear flows.~III.
  Kelvin--Helmholtz instability', ibid.\ 583--98 (1959).

\smallskip\noindent
The experimental demonstration of the Kelvin--Helmholtz instability
illustrated in the text (Fig.~117) is due to:

\item \textsc{J. R. D. Francis}, `Wave motions and the aerodynamic drag on a free oil
  surface', \emph{Phil.\ Mag.}, Ser.~7, \textbf{45}, 695--702 (1954).

\item \makebox[2em][l]{------} `Wave motions on a free oil surface',
  ibid., Ser.~8, \textbf{1}, 685--8 (1956).

\medskip
\noindent\S\,102. The distributions of $p$ and $U$ considered in this section
have been investigated by:

\item \textsc{Sir Geoffrey Taylor}, `Effect of variation in density on the stability of
  superposed streams of fluid',
  \emph{Proc.\ Roy.\ Soc.\ (London)~A}, \textbf{132}, 499--523 (1931).

\item \textsc{S. Goldstein}, `On the stability of superposed streams of fluids of different
  densities', ibid.\ 524--48 (1931).

\medskip
\noindent\S\,103. The case of an heterogeneous fluid with an exponentially
decreasing density and a linearly increasing stream velocity was first
considered by Sir Geoffrey Taylor (reference~11). The definitive discussion
of this case is due to:

\item \textsc{K. M. Case}, `Stability of an idealized atmosphere.~I. Discussion of results',
  \emph{The Physics of Fluids}, \textbf{3}, 149--54 (1960).

\item \textsc{F. J. Dyson}, `Stability of an idealized atmosphere.~II. Zeros of the
  confluent hypergeometric function', ibid.\ 155--7 (1960).

\smallskip\noindent
In connexion with Case's investigation (reference~13), see his other papers:

\item \textsc{K. M. Case}, `Taylor instability of an inverted atmosphere', ibid.\ 366--8
  (1960).

\item \makebox[2em][l]{------} `Stability of inviscid plane Couette flow', ibid.\ 143--8 (1960).

\medskip
\noindent\S\,104. The example treated in this section is due to:

\item \textsc{P. G. Drazin}, `The stability of a shear layer in an unbounded heterogeneous
  inviscid fluid', \emph{J.\ Fluid Mech.}, \textbf{4}, 214--24 (1958).

\medskip
\noindent\S\,105. See:

\item \textsc{S. Chandrasekhar}, `The effect of rotation on the development of the
  Kelvin--Helmholtz instability', unpublished.

\smallskip\noindent
For a different treatment of the same problem see:

\item \textsc{H. Solberg}, `Integrationen der atmosph\"arischen St\"orungsgleichungen, I.
  Wellenbewegungen in rotierenden, inkompressiblen Fl\"ussigkeitsschichten',
  \emph{Geofysiske Publikasjoner}, \textbf{9}, 1--120 (1928).

\smallskip\noindent
See also:

\item \textsc{V. Bjerknes}, \textsc{J. Bjerknes}, \textsc{H. Solberg},
  and \textsc{T. Bergeron}, \emph{Physikalische Hydrodynamik},
  chap.~11, Springer, Berlin, 1933.

\medskip
\noindent\S\,106. The Kelvin--Helmholtz instability in the framework of
hydromagnetics has been considered by:

\item \textsc{D. H. Michael}, `The stability of a combined current and vortex sheet in
  a perfectly conducting fluid',
  \emph{Proc.\ Camb.\ Phil.\ Soc.}, \textbf{51}, 528--32 (1955).

\item \textsc{T. G. Northrop}, `Helmholtz instability of a plasma',
  \emph{Physical Rev.}, \textbf{103}, 1150--4 (1956).

\end{enumerate}

\smallskip\noindent
Michael's discussion is restricted to a uniform density of the medium; but he
allowed for different magnetic field strengths on the two sides. Northrop's
discussion is restricted to the case when the impressed magnetic field is
transverse to the direction of streaming; but he allowed for different
strengths of the magnetic field in the two fluids.
